% !TeX encoding = UTF-8
% !TeX program = xelatex
% !TeX spellcheck = en_US

\documentclass[degree=master]{thuthesis}
  % 学位 degree:
  %   doctor | master | bachelor | postdoc
  % 学位类型 degree-type:
  %   academic（默认）| professional


% 论文基本配置，加载宏包等全局配置
% !TeX root = ./thuthesis-example.tex

% 论文基本信息配置

\thusetup{
  %******************************
  % 注意：
  %   1. 配置里面不要出现空行
  %   2. 不需要的配置信息可以删除
  %   3. 建议先阅读文档中所有关于选项的说明
  %******************************
  %
  % 输出格式
  %   选择打印版（print）或用于提交的电子版（electronic），前者会插入空白页以便直接双面打印
  %
  output = print,
  %
  % 标题
  %   可使用“\\”命令手动控制换行
  %
  title  = {清华大学学位论文 \LaTeX{} 模板\\使用示例文档 v\version},
  title* = {An Introduction to \LaTeX{} Thesis Template of Tsinghua
            University v\version},
  %
  % 学位
  %   1. 学术型
  %      - 中文
  %        需注明所属的学科门类，例如：
  %        哲学、经济学、法学、教育学、文学、历史学、理学、工学、农学、医学、
  %        军事学、管理学、艺术学
  %      - 英文
  %        博士：Doctor of Philosophy
  %        硕士：
  %          哲学、文学、历史学、法学、教育学、艺术学门类，公共管理学科
  %          填写“Master of Arts“，其它填写“Master of Science”
  %   2. 专业型
  %      直接填写专业学位的名称，例如：
  %      教育博士、工程硕士等
  %      Doctor of Education, Master of Engineering
  %   3. 本科生不需要填写
  %
  degree-name  = {工学硕士},
  degree-name* = {Master of Science},
  %
  % 培养单位
  %   填写所属院系的全名
  %
  department = {计算机科学与技术系},
  %
  % 学科
  %   1. 学术型学位
  %      获得一级学科授权的学科填写一级学科名称，其他填写二级学科名称
  %   2. 工程硕士
  %      工程领域名称
  %   3. 其他专业型学位
  %      不填写此项
  %   4. 本科生填写专业名称，第二学位论文需标注“（第二学位）”
  %
  discipline  = {计算机科学与技术},
  discipline* = {Computer Science and Technology},
  %
  % 姓名
  %
  author  = {薛瑞尼},
  author* = {Xue Ruini},
  %
  % 指导教师
  %   中文姓名和职称之间以英文逗号“,”分开，下同
  %
  supervisor  = {郑纬民, 教授},
  supervisor* = {Professor Zheng Weimin},
  %
  % 副指导教师
  %
  associate-supervisor  = {陈文光, 教授},
  associate-supervisor* = {Professor Chen Wenguang},
  %
  % 联合指导教师
  %
  % joint-supervisor  = {某某某, 教授},
  % joint-supervisor* = {Professor Mou Moumou},
  %
  % 日期
  %   使用 ISO 格式；默认为当前时间
  %
  % date = {2019-07-07},
  %
  % 是否在中文封面后的空白页生成书脊（默认 false）
  %
  include-spine = false,
  %
  % 密级和年限
  %   秘密, 机密, 绝密
  %
  % secret-level = {秘密},
  % secret-year  = {10},
  %
  % 博士后专有部分
  %
  % clc                = {分类号},
  % udc                = {UDC},
  % id                 = {编号},
  % discipline-level-1 = {计算机科学与技术},  % 流动站（一级学科）名称
  % discipline-level-2 = {系统结构},          % 专业（二级学科）名称
  % start-date         = {2011-07-01},        % 研究工作起始时间
}

%% Put any packages you would like to use here

% 表格中支持跨行
\usepackage{multirow}

% 跨页表格
\usepackage{longtable}

% 固定宽度的表格。放在 hyperref 之前的话，tabularx 里的 footnote 显示不出来。
\usepackage{tabularx}

% 表格加脚注
\usepackage{threeparttable}
\pretocmd{\TPTnoteSettings}{\footnotesize}{}{}

% 确定浮动对象的位置，可以使用 H，强制将浮动对象放到这里（可能效果很差）
\usepackage{float}

% 浮动图形控制宏包。
% 允许上一个 section 的浮动图形出现在下一个 section 的开始部分
% 该宏包提供处理浮动对象的 \FloatBarrier 命令，使所有未处
% 理的浮动图形立即被处理。这三个宏包仅供参考，未必使用：
% \usepackage[below]{placeins}
% \usepackage{floatflt} % 图文混排用宏包
% \usepackage{rotating} % 图形和表格的控制旋转

% 量和单位
\usepackage{siunitx}

% 定理类环境宏包
\usepackage{amsthm}
% 也可以使用 ntheorem
% \usepackage[amsmath,thmmarks,hyperref]{ntheorem}

% 参考文献使用 BibTeX + natbib 宏包
\usepackage[sort&compress]{natbib}
\bibliographystyle{thuthesis-numeric}      % 顺序编码制
% \bibliographystyle{thuthesis-author-year}  % 著者-出版年制
% \bibliographystyle{thuthesis-bachelor}     % 本科生参考文献的著录格式

% 参考文献使用 BibLaTeX 宏包
% \usepackage[backend=biber,style=thuthesis-numeric]{biblatex}
% \usepackage[backend=biber,style=thuthesis-author-year]{biblatex}
% \usepackage[backend=biber,style=thuthesis-bachelor]{biblatex}
% \usepackage[backend=biber,style=apa]{biblatex}
% \usepackage[backend=biber,style=mla]{biblatex}

% \addbibresource{ref/refs.bib}  % 声明 BibLaTeX 的数据库

% 给自定义的宏后面自动加空白
% \usepackage{xspace}

% 定义所有的图片文件在 figures 子目录下
\graphicspath{{figures/}}

% 定义自己常用的东西
% \def\myname{薛瑞尼}

% 数学命令
\newcommand\dif{\mathop{}\!\mathrm{d}}  % 微分符号
\newcommand\real{{\mathbf{R}}}  % 实数集
\newcommand\abs[1]{\lvert#1\rvert}
\newcommand\VECTOR{\symbf}  % 向量
\newcommand\MATRIX{\symbf}  % 矩阵
\newcommand\vn{{\VECTOR{n}}}
\newcommand\vx{{\VECTOR{x}}}
\newcommand\mA{{\MATRIX{A}}}
\newcommand\mK{{\MATRIX{K}}}

% 借用 ltxdoc 里面的几个命令方便写文档。
\DeclareRobustCommand\cs[1]{\texttt{\char`\\#1}}
\providecommand\pkg[1]{{\sffamily#1}}

% hyperref 宏包在最后调用
\usepackage{hyperref}



\begin{document}

% 封面
\maketitle

% 学位论文指导小组、公开评阅人和答辩委员会名单
% !TeX root = ../thuthesis-example.tex

\begin{committee}[name={学位论文指导小组、公开评阅人和答辩委员会名单}]

  \newcolumntype{C}[1]{@{}>{\centering\arraybackslash}p{#1}}

  \section*{指导小组名单}

  \begin{center}
    \begin{tabular}{C{3cm}C{3cm}C{9cm}@{}}
      李XX & 教授     & 清华大学 \\
      王XX & 副教授   & 清华大学 \\
      张XX & 助理教授 & 清华大学 \\
    \end{tabular}
  \end{center}


  \section*{公开评阅人名单}

  \begin{center}
    \begin{tabular}{C{3cm}C{3cm}C{9cm}@{}}
      刘XX & 教授   & 清华大学                    \\
      陈XX & 副教授 & XXXX大学                    \\
      杨XX & 研究员 & 中国XXXX科学院XXXXXXX研究所 \\
    \end{tabular}
  \end{center}


  \section*{答辩委员会名单}

  \begin{center}
    \begin{tabular}{C{2.75cm}C{2.98cm}C{4.63cm}C{4.63cm}@{}}
      主席 & 赵XX                  & 教授                    & 清华大学       \\
      委员 & 刘XX                  & 教授                    & 清华大学       \\
          & \multirow{2}{*}{杨XX} & \multirow{2}{*}{研究员} & 中国XXXX科学院 \\
          &                       &                         & XXXXXXX研究所  \\
          & 黄XX                  & 教授                    & XXXX大学       \\
          & 周XX                  & 副教授                  & XXXX大学       \\
      秘书 & 吴XX                  & 助理研究员              & 清华大学       \\
    \end{tabular}
  \end{center}

\end{committee}



% 也可以导入 Word 版转的 PDF 文件
% \begin{committee}[file=figures/committee.pdf]
% \end{committee}


% 使用授权的说明
\copyrightpage
% 将签字扫描后授权文件 scan-copyright.pdf 替换原始页面
% \copyrightpage[file=scan-copyright.pdf]

\frontmatter
% !TeX root = ../thuthesis-example.tex

% 中英文摘要和关键字

\begin{abstract}
  论文的摘要是对论文研究内容和成果的高度概括。摘要应对论文所研究的问题及其研究目
  的进行描述，对研究方法和过程进行简单介绍，对研究成果和所得结论进行概括。摘要应
  具有独立性和自明性，其内容应包含与论文全文同等量的主要信息。使读者即使不阅读全
  文，通过摘要就能了解论文的总体内容和主要成果。

  论文摘要的书写应力求精确、简明。切忌写成对论文书写内容进行提要的形式，尤其要避
  免“第 1 章……；第 2 章……；……”这种或类似的陈述方式。

  本文介绍清华大学论文模板 \thuthesis{} 的使用方法。本模板符合学校的本科、硕士、
  博士论文格式要求。

  本文的创新点主要有：
  \begin{itemize}
    \item 用例子来解释模板的使用方法；
    \item 用废话来填充无关紧要的部分；
    \item 一边学习摸索一边编写新代码。
  \end{itemize}

  关键词是为了文献标引工作、用以表示全文主要内容信息的单词或术语。关键词不超过 5
  个，每个关键词中间用分号分隔。（模板作者注：关键词分隔符不用考虑，模板会自动处
  理。英文关键词同理。）

  % 关键词用“英文逗号”分隔
  \thusetup{
    keywords = {TeX, LaTeX, CJK, 模板, 论文},
  }
\end{abstract}

\begin{abstract*}
  An abstract of a dissertation is a summary and extraction of research work
  and contributions. Included in an abstract should be description of research
  topic and research objective, brief introduction to methodology and research
  process, and summarization of conclusion and contributions of the
  research. An abstract should be characterized by independence and clarity and
  carry identical information with the dissertation. It should be such that the
  general idea and major contributions of the dissertation are conveyed without
  reading the dissertation.

  An abstract should be concise and to the point. It is a misunderstanding to
  make an abstract an outline of the dissertation and words ``the first
  chapter'', ``the second chapter'' and the like should be avoided in the
  abstract.

  Key words are terms used in a dissertation for indexing, reflecting core
  information of the dissertation. An abstract may contain a maximum of 5 key
  words, with semi-colons used in between to separate one another.
  \thusetup{
    keywords* = {TeX, LaTeX, CJK, template, thesis},
  }
\end{abstract*}


% 目录
\tableofcontents

% 插图和附表清单
\listoffigures           % 插图清单
% \listoftables            % 附表清单
% \listoffiguresandtables  % 插图和附表清单

% 符号对照表
% !TeX root = ../main.tex

\begin{denotation}[3cm]
\item[HPC] 高性能计算 (High Performance Computing)
\item[cluster] 集群
\item[Itanium] 安腾
\item[SMP] 对称多处理
\item[API] 应用程序编程接口
\item[PI] 聚酰亚胺
\item[MPI] 聚酰亚胺模型化合物，N-苯基邻苯酰亚胺
\item[PBI] 聚苯并咪唑
\item[MPBI] 聚苯并咪唑模型化合物，N-苯基苯并咪唑
\item[PY] 聚吡咙
\item[PMDA-BDA]	均苯四酸二酐与联苯四胺合成的聚吡咙薄膜
\item[$\Delta G$] 活化自由能 (Activation Free Energy)
\item[$\chi$] 传输系数 (Transmission Coefficient)
\item[$E$] 能量
\item[$m$] 质量
\item[$c$] 光速
\item[$P$] 概率
\item[$T$] 时间
\item[$v$] 速度
\item[劝学] 君子曰：学不可以已。青，取之于蓝，而青于蓝；冰，水为之，而寒于水。木
  直中绳。輮以为轮，其曲中规。虽有槁暴，不复挺者，輮使之然也。故木受绳则直，金就
  砺则利，君子博学而日参省乎己，则知明而行无过矣。吾尝终日而思矣，不如须臾之所学
  也；吾尝跂而望矣，不如登高之博见也。登高而招，臂非加长也，而见者远；顺风而呼，
  声非加疾也，而闻者彰。假舆马者，非利足也，而致千里；假舟楫者，非能水也，而绝江
  河，君子生非异也，善假于物也。积土成山，风雨兴焉；积水成渊，蛟龙生焉；积善成德，
  而神明自得，圣心备焉。故不积跬步，无以至千里；不积小流，无以成江海。骐骥一跃，
  不能十步；驽马十驾，功在不舍。锲而舍之，朽木不折；锲而不舍，金石可镂。蚓无爪牙
  之利，筋骨之强，上食埃土，下饮黄泉，用心一也。蟹六跪而二螯，非蛇鳝之穴无可寄托
  者，用心躁也。—— 荀况
\end{denotation}



% % 也可以使用 nomencl 宏包：

% \printnomenclature[3cm]

% \nomenclature{HPC}{高性能计算 (High Performance Computing)}
% \nomenclature{cluster}{集群}
% \nomenclature{Itanium}{安腾}
% \nomenclature{SMP}{对称多处理}
% \nomenclature{API}{应用程序编程接口}
% \nomenclature{PI}{聚酰亚胺}
% \nomenclature{MPI}{聚酰亚胺模型化合物，N-苯基邻苯酰亚胺}
% \nomenclature{PBI}{聚苯并咪唑}
% \nomenclature{MPBI}{聚苯并咪唑模型化合物，N-苯基苯并咪唑}
% \nomenclature{PY}{聚吡咙}
% \nomenclature{PMDA-BDA}{均苯四酸二酐与联苯四胺合成的聚吡咙薄膜}
% \nomenclature{$\Delta G$}{活化自由能 (Activation Free Energy)}
% \nomenclature{$\chi$}{传输系数 (Transmission Coefficient)}
% \nomenclature{$E$}{能量}
% \nomenclature{$m$}{质量}
% \nomenclature{$c$}{光速}
% \nomenclature{$P$}{概率}
% \nomenclature{$T$}{时间}
% \nomenclature{$v$}{速度}



% 正文部分
\mainmatter
\chapter{带 English 的标题}
\label{cha:intro}

这是 \thuthesis{} 的示例文档，基本上覆盖了模板中所有格式的设置。建议大家在使用模
板之前，除了阅读《\thuthesis{}用户手册》，这个示例文档也最好能看一看。

小老鼠偷吃热凉粉；短长虫环绕矮高粱。\footnote{韩愈（768-824），字退之，河南河阳（
  今河南孟县）人，自称郡望昌黎，世称韩昌黎。幼孤贫刻苦好学，德宗贞元八年进士。曾
  任监察御史，因上疏请免关中赋役，贬为阳山县令。后随宰相裴度平定淮西迁刑部侍郎，
  又因上表谏迎佛骨，贬潮州刺史。做过吏部侍郎，死谥文公，故世称韩吏部、韩文公。是
  唐代古文运动领袖，与柳宗元合称韩柳。诗力求险怪新奇，雄浑重气势。}


\section{封面相关}
封面的例子请参看 cover.tex。主要符号表参看 denation.tex，附录和个人简历分别参看 appendix01.tex
和 resume.tex。里面的命令都非常简单，一看即会。\footnote{你说还是看不懂？怎么会呢？}

\section{字体命令}
\label{sec:first}

苏轼（1037-1101），北宋文学家、书画家。字子瞻，号东坡居士，眉州眉山（今属四川）人
。苏洵子。嘉佑进士。神宗时曾任祠部员外郎，因反对王安石新法而求外职，任杭州通判，
知密州、徐州、湖州。后以作诗“谤讪朝廷”罪贬黄州。哲宗时任翰林学士，曾出知杭州、
颖州等，官至礼部尚书。后又贬谪惠州、儋州。北还后第二年病死常州。南宋时追谥文忠。
与父洵弟辙，合称“三苏”。在政治上属于旧党，但也有改革弊政的要求。其文汪洋恣肆，
明白畅达，为“唐宋八大家”之一。  其诗清新豪健，善用夸张比喻，在艺术表现方面独具
风格。少数诗篇也能反映民间疾苦，指责统治者的奢侈骄纵。词开豪放一派，对后代很有影
响。《念奴娇·赤壁怀古》、《水调歌头·丙辰中秋》传诵甚广。

{\kaishu 坡仙擅长行书、楷书，取法李邕、徐浩、颜真卿、杨凝式，而能自创新意。用笔丰腴
  跌宕，有天真烂漫之趣。与蔡襄、黄庭坚、米芾并称“宋四家”。能画竹，学文同，也喜
  作枯木怪石。论画主张“神似”，认为“论画以形似，见与儿童邻”；高度评价“诗中有
  画，画中有诗”的艺术造诣。诗文有《东坡七集》等。存世书迹有《答谢民师论文帖》、
  《祭黄几道文》、《前赤壁赋》、《黄州寒食诗帖》等。  画迹有《枯木怪石图》、《
  竹石图》等。}

{\fangsong 易与天地准，故能弥纶天地之道。仰以观於天文，俯以察於地理，是故知幽明之故。原
  始反终，故知死生之说。精气为物，游魂为变，是故知鬼神之情状。与天地相似，故不违。
  知周乎万物，而道济天下，故不过。旁行而不流，乐天知命，故不忧。安土敦乎仁，故
  能爱。范围天地之化而不过，曲成万物而不遗，通乎昼夜之道而知，故神无方而易无体。}

% 非本科生一般用不到幼圆与隶书字体。需要的同学可以使用 ctex-fontset-thuthesis.def 文件
{\ifcsname youyuan\endcsname\youyuan\else[无 \cs{youyuan} 字体。]\fi 有天地，然后
  万物生焉。盈天地之间者，唯万物，故受之以屯；屯者盈也，屯者物之始生也。物生必蒙，
  故受之以蒙；蒙者蒙也，物之穉也。物穉不可不养也，故受之以需；需者饮食之道也。饮
  食必有讼，故受之以讼。讼必有众起，故受之以师；师者众也。众必有所比，故受之以比；
  比者比也。比必有所畜也，故受之以小畜。物畜然后有礼，故受之以履。}

{\heiti 履而泰，然后安，故受之以泰；泰者通也。物不可以终通，故受之以否。物不可以终
  否，故受之以同人。与人同者，物必归焉，故受之以大有。有大者不可以盈，故受之以谦。
  有大而能谦，必豫，故受之以豫。豫必有随，故受之以随。以喜随人者，必有事，故受
  之以蛊；蛊者事也。}

{\ifcsname lishu\endcsname\lishu\else[无 \cs{lishu} 字体。]\fi 有事而后可大，故受
  之以临；临者大也。物大然后可观，故受之以观。可观而后有所合，故受之以噬嗑；嗑者
  合也。物不可以苟合而已，故受之以贲；贲者饰也。致饰然后亨，则尽矣，故受之以剥；
  剥者剥也。物不可以终尽，剥穷上反下，故受之以复。复则不妄矣，故受之以无妄。}

{\songti 有无妄然后可畜，故受之以大畜。物畜然后可养，故受之以颐；颐者养也。不养则不
  可动，故受之以大过。物不可以终过，故受之以坎；坎者陷也。陷必有所丽，故受之以
  离；离者丽也。}

\section{表格样本}
\label{chap1:sample:table} 

\subsection{基本表格}
\label{sec:basictable}

模板中关于表格的宏包有三个： \pkg{booktabs}、\pkg{array} 和
\pkg{longtabular}，命令有一个 \cs{hlinewd}。三线表可以用 \pkg{booktabs}
提供的 \cs{toprule}、\cs{midrule} 和 \cs{bottomrule}。它们与
\pkg{longtable} 能很好的配合使用。如果表格比较简单的话可以直接用命令
\cs{hlinewd}\marg{width} 控制。
\begin{table}[htb]
  \centering
  \begin{minipage}[t]{0.8\linewidth} % 如果想在表格中使用脚注，minipage是个不错的办法
  \caption[模板文件]{模板文件。如果表格的标题很长，那么在表格索引中就会很不美
    观，所以要像 chapter 那样在前面用中括号写一个简短的标题。这个标题会出现在索
    引中。}
  \label{tab:template-files}
    \begin{tabularx}{\linewidth}{lX}
      \toprule[1.5pt]
      {\heiti 文件名} & {\heiti 描述} \\\midrule[1pt]
      thuthesis.ins & \LaTeX{} 安装文件，docstrip\footnote{表格中的脚注} \\
      thuthesis.dtx & 所有的一切都在这里面\footnote{再来一个}。\\
      thuthesis.cls & 模板类文件。\\
      thuthesis.cfg & 模板配置文。cls 和 cfg 由前两个文件生成。\\
      thuthesis.bst    & 参考文献 Bibtex 样式文件。\\
      thuthesis.sty   & 常用的包和命令写在这里，减轻主文件的负担。\\
      \bottomrule[1.5pt]
    \end{tabularx}
  \end{minipage}
\end{table}

首先来看一个最简单的表格。表 \ref{tab:template-files} 列举了本模板主要文件及其功
能。请大家注意三线表中各条线对应的命令。这个例子还展示了如何在表格中正确使用脚注。
由于 \LaTeX{} 本身不支持在表格中使用 \cs{footnote}，所以我们不得不将表格放在
小页中，而且最好将表格的宽度设置为小页的宽度，这样脚注看起来才更美观。

\subsection{复杂表格}
\label{sec:complicatedtable}

我们经常会在表格下方标注数据来源，或者对表格里面的条目进行解释。前面的脚注是一种
不错的方法，如果你不喜欢脚注。那么完全可以在表格后面自己写注释，比如表~\ref{tab:tabexamp1}。
\begin{table}[htbp]
  \centering
  \caption{复杂表格示例 1}
  \label{tab:tabexamp1}
  \begin{minipage}[t]{0.8\textwidth} 
    \begin{tabularx}{\linewidth}{|l|X|X|X|X|}
      \hline
 \multirow{2}*{\diagbox[width=5em]{x}{y}} & \multicolumn{2}{c|}{First Half} & \multicolumn{2}{c|}{Second Half}\\\cline{2-5}
      & 1st Qtr &2nd Qtr&3rd Qtr&4th Qtr \\ \hline
      East$^{*}$ &   20.4&   27.4&   90&     20.4 \\
      West$^{**}$ &   30.6 &   38.6 &   34.6 &  31.6 \\ \hline
    \end{tabularx}\\[2pt]
    \footnotesize 注：数据来源《\thuthesis{} 使用手册》。\\
    *：东部\\
    **：西部
  \end{minipage}
\end{table}

此外，表~\ref{tab:tabexamp1} 同时还演示了另外两个功能：1）通过 \pkg{tabularx} 的
 \texttt{|X|} 扩展实现表格自动放大；2）通过命令 \cs{diagbox} 在表头部分
插入反斜线。

为了使我们的例子更接近实际情况，我会在必要的时候插入一些“无关”文字，以免太多图
表同时出现，导致排版效果不太理想。第一个出场的当然是我的最爱：风流潇洒、骏马绝尘、
健笔凌云的{\heiti 李太白}了。

李白，字太白，陇西成纪人。凉武昭王暠九世孙。或曰山东人，或曰蜀人。白少有逸才，志
气宏放，飘然有超世之心。初隐岷山，益州长史苏颋见而异之，曰：“是子天才英特，可比
相如。”天宝初，至长安，往见贺知章。知章见其文，叹曰：“子谪仙人也。”言于明皇，
召见金銮殿，奏颂一篇。帝赐食，亲为调羹，有诏供奉翰林。白犹与酒徒饮于市，帝坐沉香
亭子，意有所感，欲得白为乐章，召入，而白已醉。左右以水颒面，稍解，援笔成文，婉丽
精切。帝爱其才，数宴见。白常侍帝，醉，使高力士脱靴。力士素贵，耻之，摘其诗以激杨
贵妃。帝欲官白，妃辄沮止。白自知不为亲近所容，恳求还山。帝赐金放还。乃浪迹江湖，
终日沉饮。永王璘都督江陵，辟为僚佐。璘谋乱，兵败，白坐长流夜郎，会赦得还。族人阳
冰为当涂令，白往依之。代宗立，以左拾遗召，而白已卒。文宗时，诏以白歌诗、裴旻剑舞、
张旭草书为三绝云。集三十卷。今编诗二十五卷。\hfill —— 《全唐诗》诗人小传

浮动体的并排放置一般有两种情况：1）二者没有关系，为两个独立的浮动体；2）二者隶属
于同一个浮动体。对表格来说并排表格既可以像图~\ref{tab:parallel1}、图~\ref{tab:parallel2} 
使用小页环境，也可以如图~\ref{tab:subtable} 使用子表格来做。图的例子参见第~\ref{sec:multifig} 节。

\begin{table}[htbp]
\noindent\begin{minipage}{0.5\textwidth}
\centering
\caption{第一个并排子表格}
\label{tab:parallel1}
\begin{tabular}{p{2cm}p{2cm}}
\toprule[1.5pt]
111 & 222 \\\midrule[1pt]
222 & 333 \\\bottomrule[1.5pt]
\end{tabular}
\end{minipage}%
\begin{minipage}{0.5\textwidth}
\centering
\caption{第二个并排子表格}
\label{tab:parallel2}
\begin{tabular}{p{2cm}p{2cm}}
\toprule[1.5pt]
111 & 222 \\\midrule[1pt]
222 & 333 \\\bottomrule[1.5pt]
\end{tabular}
\end{minipage}
\end{table}

然后就是忧国忧民，诗家楷模杜工部了。杜甫，字子美，其先襄阳人，曾祖依艺为巩令，因
居巩。甫天宝初应进士，不第。后献《三大礼赋》，明皇奇之，召试文章，授京兆府兵曹参
军。安禄山陷京师，肃宗即位灵武，甫自贼中遁赴行在，拜左拾遗。以论救房琯，出为华州
司功参军。关辅饥乱，寓居同州同谷县，身自负薪采梠，餔糒不给。久之，召补京兆府功曹，
道阻不赴。严武镇成都，奏为参谋、检校工部员外郎，赐绯。武与甫世旧，待遇甚厚。乃于
成都浣花里种竹植树，枕江结庐，纵酒啸歌其中。武卒，甫无所依，乃之东蜀就高適。既至
而適卒。是岁，蜀帅相攻杀，蜀大扰。甫携家避乱荆楚，扁舟下峡，未维舟而江陵亦乱。乃
溯沿湘流，游衡山，寓居耒阳。卒年五十九。元和中，归葬偃师首阳山，元稹志其墓。天宝
间，甫与李白齐名，时称李杜。然元稹之言曰：“李白壮浪纵恣，摆去拘束，诚亦差肩子美
矣。至若铺陈终始，排比声韵，大或千言，次犹数百，词气豪迈，而风调清深，属对律切，
而脱弃凡近，则李尚不能历其藩翰，况堂奥乎。”白居易亦云：“杜诗贯穿古今，  尽工尽
善，殆过于李。”元、白之论如此。盖其出处劳佚，喜乐悲愤，好贤恶恶，一见之于诗。而
又以忠君忧国、伤时念乱为本旨。读其诗可以知其世，故当时谓之“诗史”。旧集诗文共六
十卷，今编诗十九卷。

\begin{table}[htbp]
\centering
\caption{并排子表格}
\label{tab:subtable}
\subcaptionbox{第一个子表格}
{
\begin{tabular}{p{2cm}p{2cm}}
\toprule[1.5pt]
111 & 222 \\\midrule[1pt]
222 & 333 \\\bottomrule[1.5pt]
\end{tabular}
}
\hskip2cm
\subcaptionbox{第二个子表格}
{
\begin{tabular}{p{2cm}p{2cm}}
\toprule[1.5pt]
111 & 222 \\\midrule[1pt]
222 & 333 \\\bottomrule[1.5pt]
\end{tabular}
}
\end{table}

不可否认 \LaTeX{} 的表格功能没有想象中的那么强大，不过只要你足够认真，足够细致，那么
同样可以排出来非常复杂非常漂亮的表格。请参看表~\ref{tab:tabexamp2}。
\begin{table}[htbp]
  \centering\dawu[1.3]
  \caption{复杂表格示例 2}
  \label{tab:tabexamp2}
  \begin{tabular}[c]{|m{1.5cm}|c|c|c|c|c|c|}\hline
    \multicolumn{2}{|c|}{Network Topology} & \# of nodes & 
    \multicolumn{3}{c|}{\# of clients} & Server \\\hline
    GT-ITM & Waxman Transit-Stub & 600 &
    \multirow{2}{2em}{2\%}& 
    \multirow{2}{2em}{10\%}& 
    \multirow{2}{2em}{50\%}& 
    \multirow{2}{1.2in}{Max. Connectivity}\\\cline{1-3}
    \multicolumn{2}{|c|}{Inet-2.1} & 6000 & & & &\\\hline
    \multirow{2}{1.5cm}{Xue} & Rui  & Ni &\multicolumn{4}{c|}{\multirow{2}*{\thuthesis}}\\\cline{2-3}
    & \multicolumn{2}{c|}{ABCDEF} &\multicolumn{4}{c|}{} \\\hline
\end{tabular}
\end{table}

最后就是清新飘逸、文约意赅、空谷绝响的王大侠了。王维，字摩诘，河东人。工书画，与
弟缙俱有俊才。开元九年，进士擢第，调太乐丞。坐累为济州司仓参军，历右拾遗、监察御
史、左补阙、库部郎中，拜吏部郎中。天宝末，为给事中。安禄山陷两都，维为贼所得，服
药阳喑，拘于菩提寺。禄山宴凝碧池，维潜赋诗悲悼，闻于行在。贼平，陷贼官三等定罪，
特原之，责授太子中允，迁中庶子、中书舍人。复拜给事中，转尚书右丞。维以诗名盛于开
元、天宝间，宁薛诸王驸马豪贵之门，无不拂席迎之。得宋之问辋川别墅，山水绝胜，与道
友裴迪，浮舟往来，弹琴赋诗，啸咏终日。笃于奉佛，晚年长斋禅诵。一日，忽索笔作书
数纸，别弟缙及平生亲故，舍笔而卒。赠秘书监。宝应中，代宗问缙：“朕常于诸王坐闻维
乐章，今存几何？”缙集诗六卷，文四卷，表上之。敕答云，卿伯氏位列先朝，名高希代。
抗行周雅，长揖楚辞。诗家者流，时论归美。克成编录，叹息良深。殷璠谓维诗词秀调雅，
意新理惬。在泉成珠，著壁成绘。苏轼亦云：“维诗中有画，画中有诗也。”今编诗四卷。

要想用好论文模板还是得提前学习一些 \TeX/\LaTeX{}的相关知识，具备一些基本能力，掌
握一些常见技巧，否则一旦遇到问题还真是比较麻烦。我们见过很多这样的同学，一直以来
都是使用 Word 等字处理工具，以为 \LaTeX{}模板的用法也应该类似，所以就沿袭同样的思
路来对待这种所见非所得的排版工具，结果被折腾的焦头烂额，疲惫不堪。

如果您要排版的表格长度超过一页，那么推荐使用 \pkg{longtable} 或者 \pkg{supertabular}
宏包，模板对 \pkg{longtable} 进行了相应的设置，所以用起来可能简单一些。
表~\ref{tab:performance} 就是 \pkg{longtable} 的简单示例。
\begin{longtable}[c]{c*{6}{r}}
\caption{实验数据}\label{tab:performance}\\
\toprule[1.5pt]
 测试程序 & \multicolumn{1}{c}{正常运行} & \multicolumn{1}{c}{同步} & \multicolumn{1}{c}{检查点} & \multicolumn{1}{c}{卷回恢复}
& \multicolumn{1}{c}{进程迁移} & \multicolumn{1}{c}{检查点} \\
& \multicolumn{1}{c}{时间 (s)}& \multicolumn{1}{c}{时间 (s)}&
\multicolumn{1}{c}{时间 (s)}& \multicolumn{1}{c}{时间 (s)}& \multicolumn{1}{c}{
  时间 (s)}&  文件（KB）\\\midrule[1pt]
\endfirsthead
\multicolumn{7}{c}{续表~\thetable\hskip1em 实验数据}\\
\toprule[1.5pt]
 测试程序 & \multicolumn{1}{c}{正常运行} & \multicolumn{1}{c}{同步} & \multicolumn{1}{c}{检查点} & \multicolumn{1}{c}{卷回恢复}
& \multicolumn{1}{c}{进程迁移} & \multicolumn{1}{c}{检查点} \\
& \multicolumn{1}{c}{时间 (s)}& \multicolumn{1}{c}{时间 (s)}&
\multicolumn{1}{c}{时间 (s)}& \multicolumn{1}{c}{时间 (s)}& \multicolumn{1}{c}{
  时间 (s)}&  文件（KB）\\\midrule[1pt]
\endhead
\hline
\multicolumn{7}{r}{续下页}
\endfoot
\endlastfoot
CG.A.2 & 23.05 & 0.002 & 0.116 & 0.035 & 0.589 & 32491 \\
CG.A.4 & 15.06 & 0.003 & 0.067 & 0.021 & 0.351 & 18211 \\
CG.A.8 & 13.38 & 0.004 & 0.072 & 0.023 & 0.210 & 9890 \\
CG.B.2 & 867.45 & 0.002 & 0.864 & 0.232 & 3.256 & 228562 \\
CG.B.4 & 501.61 & 0.003 & 0.438 & 0.136 & 2.075 & 123862 \\
CG.B.8 & 384.65 & 0.004 & 0.457 & 0.108 & 1.235 & 63777 \\
MG.A.2 & 112.27 & 0.002 & 0.846 & 0.237 & 3.930 & 236473 \\
MG.A.4 & 59.84 & 0.003 & 0.442 & 0.128 & 2.070 & 123875 \\
MG.A.8 & 31.38 & 0.003 & 0.476 & 0.114 & 1.041 & 60627 \\
MG.B.2 & 526.28 & 0.002 & 0.821 & 0.238 & 4.176 & 236635 \\
MG.B.4 & 280.11 & 0.003 & 0.432 & 0.130 & 1.706 & 123793 \\
MG.B.8 & 148.29 & 0.003 & 0.442 & 0.116 & 0.893 & 60600 \\
LU.A.2 & 2116.54 & 0.002 & 0.110 & 0.030 & 0.532 & 28754 \\
LU.A.4 & 1102.50 & 0.002 & 0.069 & 0.017 & 0.255 & 14915 \\
LU.A.8 & 574.47 & 0.003 & 0.067 & 0.016 & 0.192 & 8655 \\
LU.B.2 & 9712.87 & 0.002 & 0.357 & 0.104 & 1.734 & 101975 \\
LU.B.4 & 4757.80 & 0.003 & 0.190 & 0.056 & 0.808 & 53522 \\
LU.B.8 & 2444.05 & 0.004 & 0.222 & 0.057 & 0.548 & 30134 \\
EP.A.2 & 123.81 & 0.002 & 0.010 & 0.003 & 0.074 & 1834 \\
EP.A.4 & 61.92 & 0.003 & 0.011 & 0.004 & 0.073 & 1743 \\
EP.A.8 & 31.06 & 0.004 & 0.017 & 0.005 & 0.073 & 1661 \\
EP.B.2 & 495.49 & 0.001 & 0.009 & 0.003 & 0.196 & 2011 \\
EP.B.4 & 247.69 & 0.002 & 0.012 & 0.004 & 0.122 & 1663 \\
EP.B.8 & 126.74 & 0.003 & 0.017 & 0.005 & 0.083 & 1656 \\
\bottomrule[1.5pt]
\end{longtable}

\subsection{其它}
\label{sec:tableother}
如果不想让某个表格或者图片出现在索引里面，请使用命令 \cs{caption*}。
这个命令不会给表格编号，也就是出来的只有标题文字而没有“表~XX”，“图~XX”，否则
索引里面序号不连续就显得不伦不类，这也是 \LaTeX{} 里星号命令默认的规则。

有这种需求的多是本科同学的英文资料翻译部分，如果你觉得附录中英文原文中的表格和图
片显示成“  表”和“图”很不协调的话，一个很好的办法就是用 \cs{caption*}，参数
随便自己写，比如不守规矩的表~1.111 和图~1.111 能满足这种特殊需要（可以参看附录部
分）。
\begin{table}[ht]
  \begin{minipage}{0.4\linewidth}
    \centering
    \caption*{表~1.111\quad 这是一个手动编号，不出现在索引中的表格。}
    \label{tab:badtabular}
      \framebox(150,50)[c]{\thuthesis}
  \end{minipage}%
  \hfill%
  \begin{minipage}{0.4\linewidth}
    \centering
    \caption*{Figure~1.111\quad 这是一个手动编号，不出现在索引中的图。}
    \label{tab:badfigure}
    \framebox(150,50)[c]{薛瑞尼}
  \end{minipage}
\end{table}

如果你的确想让它编号，但又不想让它出现在索引中的话，那就自己看看代码改一改吧，我
目前不打算给模板增加这种另类命令。

最后，虽然大家不一定会独立使用小页，但是关于小页中的脚注还是有必要提一下。请看下
面的例子。

\begin{minipage}[t]{\linewidth-2\parindent}
  柳宗元，字子厚（773-819），河东（今永济县）人\footnote{山西永济水饺。}，是唐代
  杰出的文学家，哲学家，同时也是一位政治改革家。与韩愈共同倡导唐代古文运动，并称
  韩柳\footnote{唐宋八大家之首二位。}。
\end{minipage}

唐朝安史之乱后，宦官专权，藩镇割据，土地兼并日渐严重，社会生产破坏严重，民不聊生。柳宗
元对这种社会现实极为不满，他积极参加了王叔文领导的“永济革新”，并成为这一
运动的中坚人物。他们革除弊政，打击权奸，触犯了宦官和官僚贵族利益，在他们的联合反
扑下，改革失败了，柳宗元被贬为永州司马。

\section{定理环境}
\label{sec:theorem}

给大家演示一下各种和证明有关的环境：

\begin{assumption}
待月西厢下，迎风户半开；隔墙花影动，疑是玉人来。
\begin{eqnarray}
  \label{eq:eqnxmp}
  c & = & a^2 - b^2\\
    & = & (a+b)(a-b)
\end{eqnarray}
\end{assumption}

千辛万苦，历尽艰难，得有今日。然相从数千里，未曾哀戚。今将渡江，方图百年欢笑，如
何反起悲伤？（引自《杜十娘怒沉百宝箱》）

\begin{definition}
子曰：「道千乘之国，敬事而信，节用而爱人，使民以时。」
\end{definition}

千古第一定义！问世间、情为何物，只教生死相许？天南地北双飞客，老翅几回寒暑。欢乐趣，离别苦，就中更有痴儿女。
君应有语，渺万里层云，千山暮雪，只影向谁去？

横汾路，寂寞当年箫鼓，荒烟依旧平楚。招魂楚些何嗟及，山鬼暗谛风雨。天也妒，未信与，莺儿燕子俱黄土。
千秋万古，为留待骚人，狂歌痛饮，来访雁丘处。

\begin{proposition}
 曾子曰：「吾日三省吾身 —— 为人谋而不忠乎？与朋友交而不信乎？传不习乎？」
\end{proposition}

多么凄美的命题啊！其日牛马嘶，新妇入青庐，奄奄黄昏后，寂寂人定初，我命绝今日，
魂去尸长留，揽裙脱丝履，举身赴清池，府吏闻此事，心知长别离，徘徊庭树下，自挂东南
枝。

\begin{remark}
天不言自高，水不言自流。
\begin{gather*}
\begin{split} 
\varphi(x,z)
&=z-\gamma_{10}x-\gamma_{mn}x^mz^n\\
&=z-Mr^{-1}x-Mr^{-(m+n)}x^mz^n
\end{split}\\[6pt]
\begin{align} \zeta^0&=(\xi^0)^2,\\
\zeta^1 &=\xi^0\xi^1,\\
\zeta^2 &=(\xi^1)^2,
\end{align}
\end{gather*}
\end{remark}

天尊地卑，乾坤定矣。卑高以陈，贵贱位矣。 动静有常，刚柔断矣。方以类聚，物以群分，
吉凶生矣。在天成象，在地成形，变化见矣。鼓之以雷霆，润之以风雨，日月运行，一寒一
暑，乾道成男，坤道成女。乾知大始，坤作成物。乾以易知，坤以简能。易则易知，简则易
从。易知则有亲，易从则有功。有亲则可久，有功则可大。可久则贤人之德，可大则贤人之
业。易简，而天下矣之理矣；天下之理得，而成位乎其中矣。

\begin{axiom}
两点间直线段距离最短。  
\begin{align}
x&\equiv y+1\pmod{m^2}\\
x&\equiv y+1\mod{m^2}\\
x&\equiv y+1\pod{m^2}
\end{align}
\end{axiom}

《彖曰》：大哉乾元，万物资始，乃统天。云行雨施，品物流形。大明始终，六位时成，时
乘六龙以御天。乾道变化，各正性命，保合大和，乃利贞。首出庶物，万国咸宁。

《象曰》：天行健，君子以自强不息。潜龙勿用，阳在下也。见龙再田，德施普也。终日乾
乾，反复道也。或跃在渊，进无咎也。飞龙在天，大人造也。亢龙有悔，盈不可久也。用九，
天德不可为首也。 　　

\begin{lemma}
《猫和老鼠》是我最爱看的动画片。
\begin{multline*}%\tag*{[a]} % 这个不出现在索引中
\int_a^b\biggl\{\int_a^b[f(x)^2g(y)^2+f(y)^2g(x)^2]
 -2f(x)g(x)f(y)g(y)\,dx\biggr\}\,dy \\
 =\int_a^b\biggl\{g(y)^2\int_a^bf^2+f(y)^2
  \int_a^b g^2-2f(y)g(y)\int_a^b fg\biggr\}\,dy
\end{multline*}
\end{lemma}

行行重行行，与君生别离。相去万余里，各在天一涯。道路阻且长，会面安可知。胡马依北
风，越鸟巢南枝。相去日已远，衣带日已缓。浮云蔽白日，游子不顾返。思君令人老，岁月
忽已晚。  弃捐勿复道，努力加餐饭。

\begin{theorem}\label{the:theorem1}
犯我强汉者，虽远必诛\hfill —— 陈汤（汉）
\end{theorem}
\begin{subequations}
\begin{align}
y & = 1 \\
y & = 0
\end{align}
\end{subequations}
道可道，非常道。名可名，非常名。无名天地之始；有名万物之母。故常无，欲以观其妙；
常有，欲以观其徼。此两者，同出而异名，同谓之玄。玄之又玄，众妙之门。上善若水。水
善利万物而不争，处众人之所恶，故几于道。曲则全，枉则直，洼则盈，敝则新，少则多，
多则惑。人法地，地法天，天法道，道法自然。知人者智，自知者明。胜人者有力，自胜
者强。知足者富。强行者有志。不失其所者久。死而不亡者寿。

\begin{proof}
燕赵古称多感慨悲歌之士。董生举进士，连不得志于有司，怀抱利器，郁郁适兹土，吾
知其必有合也。董生勉乎哉？

夫以子之不遇时，苟慕义强仁者，皆爱惜焉，矧燕、赵之士出乎其性者哉！然吾尝闻
风俗与化移易，吾恶知其今不异于古所云邪？聊以吾子之行卜之也。董生勉乎哉？

吾因子有所感矣。为我吊望诸君之墓，而观于其市，复有昔时屠狗者乎？为我谢
曰：“明天子在上，可以出而仕矣！” \hfill —— 韩愈《送董邵南序》
\end{proof}

\begin{corollary}
  四川话配音的《猫和老鼠》是世界上最好看最好听最有趣的动画片。
\begin{alignat}{3}
V_i & =v_i - q_i v_j, & \qquad X_i & = x_i - q_i x_j,
 & \qquad U_i & = u_i,
 \qquad \text{for $i\ne j$;}\label{eq:B}\\
V_j & = v_j, & \qquad X_j & = x_j,
  & \qquad U_j & u_j + \sum_{i\ne j} q_i u_i.
\end{alignat}
\end{corollary}

迢迢牵牛星，皎皎河汉女。
纤纤擢素手，札札弄机杼。
终日不成章，泣涕零如雨。
河汉清且浅，相去复几许。
盈盈一水间，脉脉不得语。

\begin{example}
  大家来看这个例子。
\begin{equation}
\label{ktc}
\left\{\begin{array}{l}
\nabla f({\mbox{\boldmath $x$}}^*)-\sum\limits_{j=1}^p\lambda_j\nabla g_j({\mbox{\boldmath $x$}}^*)=0\\[0.3cm]
\lambda_jg_j({\mbox{\boldmath $x$}}^*)=0,\quad j=1,2,\cdots,p\\[0.2cm]
\lambda_j\ge 0,\quad j=1,2,\cdots,p.
\end{array}\right.
\end{equation}
\end{example}

\begin{exercise}
  清列出 Andrew S. Tanenbaum 和 W. Richard Stevens 的所有著作。
\end{exercise}

\begin{conjecture} \textit{Poincare Conjecture} If in a closed three-dimensional
  space, any closed curves can shrink to a point continuously, this space can be
  deformed to a sphere.
\end{conjecture}

\begin{problem}
 回答还是不回答，是个问题。 
\end{problem}

如何引用定理~\ref{the:theorem1} 呢？加上 \cs{label} 使用 \cs{ref} 即可。妾发
初覆额，折花门前剧。郎骑竹马来，绕床弄青梅。同居长干里，两小无嫌猜。 十四为君妇，
羞颜未尝开。低头向暗壁，千唤不一回。十五始展眉，愿同尘与灰。常存抱柱信，岂上望夫
台。 十六君远行，瞿塘滟滪堆。五月不可触，猿声天上哀。门前迟行迹，一一生绿苔。苔深
不能扫，落叶秋风早。八月蝴蝶来，双飞西园草。感此伤妾心，坐愁红颜老。

\section{参考文献}
\label{sec:bib}
当然参考文献可以直接写 bibitem，虽然费点功夫，但是好控制，各种格式可以自己随意改
写。

本模板推荐使用 BIB\TeX，样式文件为 \texttt{thuthesis.bst}，基本符合学校的参考文献格
式（如专利等引用未加详细测试）。看看这个例子，关于书的~\cite{tex, companion,
  ColdSources}，还有这些~\cite{Krasnogor2004e, clzs, zjsw}，关于杂志
的~\cite{ELIDRISSI94, MELLINGER96, SHELL02}，硕士论文~\cite{zhubajie,
  metamori2004}，博士论文~\cite{shaheshang, FistSystem01}，标准文
件~\cite{IEEE-1363}，会议论文~\cite{DPMG,kocher99}，技术报告~\cite{NPB2}，电子文
献~\cite{chuban2001,oclc2000}。中文参考文献~\cite{cnarticle}应增
加 \texttt{lang=``zh''} 字段，以便进行相应处理。另外，本模板对中文文
献~\cite{cnproceed}的支持并不是十全十美，如果有不如意的地方，请手动修
改 \texttt{bbl} 文件。

有时候不想要上标，那么可以这样~\inlinecite{shaheshang}，这个非常重要。

有时候一些参考文献没有纸质出处，需要标注 URL。缺省情况下，URL 不会在连字符处断行，
这可能使得用连字符代替空格的网址分行很难看。如果需要，可以将模板类文件中
\begin{verbatim}
\RequirePackage{hyperref}
\end{verbatim}
一行改为：
\begin{verbatim}
\PassOptionsToPackage{hyphens}{url}
\RequirePackage{hyperref}
\end{verbatim}
使得连字符处可以断行。更多设置可以参考 \texttt{url} 宏包文档。

\section{公式}
\label{sec:equation}
贝叶斯公式如式~(\ref{equ:chap1:bayes})，其中 $p(y|\mathbf{x})$ 为后验；
$p(\mathbf{x})$ 为先验；分母 $p(\mathbf{x})$ 为归一化因子。
\begin{equation}
\label{equ:chap1:bayes}
p(y|\mathbf{x}) = \frac{p(\mathbf{x},y)}{p(\mathbf{x})}=
\frac{p(\mathbf{x}|y)p(y)}{p(\mathbf{x})} 
\end{equation}

论文里面公式越多，\TeX{} 就越 happy。再看一个 \pkg{amsmath} 的例子：
\newcommand{\envert}[1]{\left\lvert#1\right\rvert} 
\begin{equation}\label{detK2}
\det\mathbf{K}(t=1,t_1,\dots,t_n)=\sum_{I\in\mathbf{n}}(-1)^{\envert{I}}
\prod_{i\in I}t_i\prod_{j\in I}(D_j+\lambda_jt_j)\det\mathbf{A}
^{(\lambda)}(\overline{I}|\overline{I})=0.
\end{equation} 

前面定理示例部分列举了很多公式环境，可以说把常见的情况都覆盖了，大家在写公式的时
候一定要好好看 \pkg{amsmath} 的文档，并参考模板中的用法：
\begin{multline*}%\tag{[b]} % 这个出现在索引中的
\int_a^b\biggl\{\int_a^b[f(x)^2g(y)^2+f(y)^2g(x)^2]
 -2f(x)g(x)f(y)g(y)\,dx\biggr\}\,dy \\
 =\int_a^b\biggl\{g(y)^2\int_a^bf^2+f(y)^2
  \int_a^b g^2-2f(y)g(y)\int_a^b fg\biggr\}\,dy
\end{multline*}

其实还可以看看这个多级规划：
\begin{equation}\label{bilevel}
\left\{\begin{array}{l}
\max\limits_{{\mbox{\footnotesize\boldmath $x$}}} F(x,y_1^*,y_2^*,\cdots,y_m^*)\\[0.2cm]
\mbox{subject to:}\\[0.1cm]
\qquad G(x)\le 0\\[0.1cm]
\qquad(y_1^*,y_2^*,\cdots,y_m^*)\mbox{ solves problems }(i=1,2,\cdots,m)\\[0.1cm]
\qquad\left\{\begin{array}{l}
    \max\limits_{{\mbox{\footnotesize\boldmath $y_i$}}}f_i(x,y_1,y_2,\cdots,y_m)\\[0.2cm]
    \mbox{subject to:}\\[0.1cm]
    \qquad g_i(x,y_1,y_2,\cdots,y_m)\le 0.
    \end{array}\right.
\end{array}\right.
\end{equation}
这些跟规划相关的公式都来自于刘宝碇老师《不确定规划》的课件。


%%% Local Variables: 
%%% mode: latex
%%% TeX-master: t
%%% End: 

\chapter{中华人民共和国}
\label{cha:china}

\section{其它例子}
\label{sec:other}

在第~\ref{cha:intro} 章中我们学习了贝叶斯公式~(\ref{equ:chap1:bayes})，这里我们复
习一下：
\begin{equation}
\label{equ:chap2:bayes}
p(y|\mathbf{x}) = \frac{p(\mathbf{x},y)}{p(\mathbf{x})}=
\frac{p(\mathbf{x}|y)p(y)}{p(\mathbf{x})} 
\end{equation}

\subsection{绘图}
\label{sec:draw}

本模板不再预先装载任何绘图包（如 \textsf{pstricks，pgf} 等），完全由你自己来决定。
个人觉得 \textsf{pgf} 不错，不依赖于 Postscript。此外还有很多针对 \LaTeX{} 的
 GUI 作图工具，如 XFig(jFig), WinFig, Tpx, Ipe, Dia, Inkscape, LaTeXPiX,
jPicEdt, jaxdraw 等等。

\subsection{插图}
\label{sec:graphs}

强烈推荐《\LaTeXe 插图指南》！关于子图形的使用细节请参看 \textsf{subcaption} 宏包的说明文档。

\subsubsection{一个图形}
\label{sec:onefig}
一般图形都是处在浮动环境中。之所以称为浮动是指最终排版效果图形的位置不一定与源文
件中的位置对应\footnote{This is not a bug, but a feature of \LaTeX!}，这也是刚使
用 \LaTeX{} 同学可能遇到的问题。如果要强制固定浮动图形的位置，请使用 \textsf{float} 宏包，
它提供了 \texttt{[H]} 参数，比如图~\ref{fig:xfig1}。
\begin{figure}[H] % use float package if you want it here
  \centering
  \includegraphics{hello}
  \caption{利用 Xfig 制图}
  \label{fig:xfig1}
\end{figure}

大学之道，在明明德，在亲民，在止于至善。知止而后有定；定而后能静；静而后能安；安
而后能虑；虑而后能得。物有本末，事有终始。知所先后，则近道矣。古之欲明明德于天
下者，先治其国；欲治其国者，先齐其家；欲齐其家者，先修其身；欲修其身者，先正其心；
欲正其心者，先诚其意；欲诚其意者，先致其知；致知在格物。物格而后知至；知至而后
意诚；意诚而后心正；心正而后身 修；身修而后家齐；家齐而后国治；国治而后天下
平。自天子以至于庶人，壹是皆以修身为本。其本乱而未治者 否矣。其所厚者薄，而其所
薄者厚，未之有也！

\hfill \pozhehao《大学》


\subsubsection{多个图形}
\label{sec:multifig}

如果多个图形相互独立，并不共用一个图形计数器，那么用 \verb|minipage| 或者
\verb|parbox| 就可以。否则，请参看图~\ref{fig:big1-subcaptionbox}，它包含两个小图，分别是图~\ref{fig:subfig1} 
和图~\ref{fig:subfig2}。推荐使用\verb|\subcaptionbox|，
因为可以像图~\ref{fig:big1-subcaptionbox} 那样对齐子图的标题，
也可以使用\textsf{subcaption}宏包的\verb|\subcaption|（放在minipage中，用法同\verb|\caption|）
或是 subfigure 、 subtable环境，像图~\ref{fig:big1-subfigure}，不要再用 \verb|\subfloat|、
\verb|\subfigure| 和 \verb|\subtable|。
\begin{figure}[h]
  \centering%
  \subcaptionbox{第一个小图形\label{fig:subfig1}}
  [3cm] %标题的长度，超过则会换行，如下一个小图。
    {\includegraphics[height=3cm]{thu-fig-logo}}
      \hspace{4em}%
  \subcaptionbox{第二个小图形，注意这个图略矮些。如果标题很长的话，它会自动换行\label{fig:subfig2}}
      {\includegraphics[height=2cm]{thu-text-logo}}
  \caption{包含子图形的大图形(subcaptionbox示例)}
  \label{fig:big1-subcaptionbox}
\end{figure}
\begin{figure}[h]
  \centering%
  \begin{subfigure}{3cm}
    \includegraphics[height=3cm]{thu-fig-logo}
    \caption{第一个小图形}
  \end{subfigure}
  \hspace{4em}%
  \begin{subfigure}{0.5\textwidth}
    \includegraphics[height=2cm]{thu-text-logo}
    \caption{第二个小图形，注意这个图略矮些。subfigure中同一行的子图在顶端对齐。}
  \end{subfigure}
  \caption{包含子图形的大图形(subfigure示例)}
  \label{fig:big1-subfigure}
\end{figure}
古之学者必有师。师者，所以传道受业解惑也。人非生而知之者，孰能无惑？惑而不从师，
其为惑也，终不解矣。生乎吾前，其闻道也固先乎吾，吾从而师之；生乎吾後，其闻道也亦
先乎吾，吾从而师之。吾师道也，夫庸知其年之先後生於吾乎！是故无贵无贱无长无少，道
之所存，师之所存也。

嗟乎！师道之不传也久矣，欲人之无惑也难矣。古之圣人，其出人也远矣，犹且从师而问焉；
今之众人，其下圣人也亦远矣，而耻学於师。是故圣益圣，愚益愚。圣人之所以为圣，愚
人之所以为愚，其皆出於此乎？爱其子，择师而教之，於其身也，则耻师焉，惑焉。彼童子
之师，授之书而习其句读者，非吾所谓传其道、解其惑者也。句读之不知，惑之不解，或师
焉，或不焉，小学而大遗，吾未见其明也。巫医、乐师、百工之人不耻相师，  士大夫之族
曰“师”曰“弟子”之云者，则群聚而笑之。问之，则曰：彼与彼年相若也，道相似也，位
卑则足羞，官盛则近谀。呜呼！师道之不复，可知矣。巫医、乐师、百工之人。吾子不齿，
今其智乃反不能及，其可怪也欤！圣人无常师。孔子师郯子、苌子、师襄、老聃。郯子之徒，
其贤不及孔子。孔子曰：“三人行，必有我师。”是故弟子不必不如师，师不必贤於弟子。
闻道有先後，术业有专攻，如是而已。

如果要把编号的两个图形并排，那么小页就非常有用了：
\begin{figure}
\begin{minipage}{0.48\textwidth}
  \centering
  \includegraphics[height=2cm]{thu-whole-logo}
  \caption{并排第一个图}
  \label{fig:parallel1}
\end{minipage}\hfill
\begin{minipage}{0.48\textwidth}
  \centering
  \includegraphics[height=2cm]{thu-whole-logo}
  \caption{并排第二个图}
  \label{fig:parallel2}
\end{minipage}
\end{figure}

李氏子蟠，年十七，好古文、六艺，经传皆通习之，不拘於时，学於余。余嘉其能行古
道，作师说以贻之。

\hfill \pozhehao 韩愈（唐）




% 其他部分
\backmatter

% 参考文献
\bibliography{ref/refs}  % 参考文献使用 BibTeX 编译
% \printbibliography       % 参考文献使用 BibLaTeX 编译

% 附录
\appendix
\chapter{单目标规划}

\section{单目标规划}
北冥有鱼，其名为鲲。鲲之大，不知其几千里也。化而为鸟，其名为鹏。鹏之背，不知其几
千里也。怒而飞，其翼若垂天之云。是鸟也，海运则将徙于南冥。南冥者，天池也。
\begin{equation}\tag*{(123)}
 p(y|\mathbf{x}) = \frac{p(\mathbf{x},y)}{p(\mathbf{x})}=
\frac{p(\mathbf{x}|y)p(y)}{p(\mathbf{x})}
\end{equation}

吾生也有涯，而知也无涯。以有涯随无涯，殆已！已而为知者，殆而已矣！为善无近名，为
恶无近刑，缘督以为经，可以保身，可以全生，可以养亲，可以尽年。

\subsection{线性规划}
庖丁为文惠君解牛，手之所触，肩之所倚，足之所履，膝之所倚，砉然响然，奏刀騞然，莫
不中音，合于桑林之舞，乃中经首之会。
\begin{table}[ht]
\centering
  \centering
  \caption{这是手动编号但不出现在索引中的一个表格例子}
  \label{tab:badtabular3}
  \begin{tabular}[c]{|m{1.5cm}|c|c|c|c|c|c|}\hline
    \multicolumn{2}{|c|}{Network Topology} & \# of nodes &
    \multicolumn{3}{c|}{\# of clients} & Server \\\hline
    GT-ITM & Waxman Transit-Stub & 600 &
    \multirow{2}{2em}{2\%}&
    \multirow{2}{2em}{10\%}&
    \multirow{2}{2em}{50\%}&
    \multirow{2}{1.2in}{Max. Connectivity}\\\cline{1-3}
    \multicolumn{2}{|c|}{Inet-2.1} & 6000 & & & &\\\hline
    \multirow{2}{1.5cm}{Xue} & Rui  & Ni &\multicolumn{4}{c|}{\multirow{2}*{\thuthesis}}\\\cline{2-3}
    & \multicolumn{2}{c|}{ABCDEF} &\multicolumn{4}{c|}{} \\\hline
\end{tabular}
\end{table}

文惠君曰：“嘻，善哉！技盖至此乎？”庖丁释刀对曰：“臣之所好者道也，进乎技矣。始臣之
解牛之时，所见无非全牛者；三年之后，未尝见全牛也；方今之时，臣以神遇而不以目视，
官知止而神欲行。依乎天理，批大郤，导大窾，因其固然。技经肯綮之未尝，而况大坬乎！
良庖岁更刀，割也；族庖月更刀，折也；今臣之刀十九年矣，所解数千牛矣，而刀刃若新发
于硎。彼节者有间而刀刃者无厚，以无厚入有间，恢恢乎其于游刃必有余地矣。是以十九年
而刀刃若新发于硎。虽然，每至于族，吾见其难为，怵然为戒，视为止，行为迟，动刀甚微，
謋然已解，如土委地。提刀而立，为之而四顾，为之踌躇满志，善刀而藏之。”

文惠君曰：“善哉！吾闻庖丁之言，得养生焉。”


\subsection{非线性规划}
孔子与柳下季为友，柳下季之弟名曰盗跖。盗跖从卒九千人，横行天下，侵暴诸侯。穴室枢
户，驱人牛马，取人妇女。贪得忘亲，不顾父母兄弟，不祭先祖。所过之邑，大国守城，小
国入保，万民苦之。孔子谓柳下季曰：“夫为人父者，必能诏其子；为人兄者，必能教其弟。
若父不能诏其子，兄不能教其弟，则无贵父子兄弟之亲矣。今先生，世之才士也，弟为盗
跖，为天下害，而弗能教也，丘窃为先生羞之。丘请为先生往说之。”
\begin{figure}[h]
  \centering
  \includegraphics{thu-whole-logo.pdf}
  \caption{这是手动编号但不出现索引中的图片的例子}
  \label{tab:badfigure3}
\end{figure}

柳下季曰：“先生言为人父者必能诏其子，为人兄者必能教其弟，若子不听父之诏，弟不受
兄之教，虽今先生之辩，将奈之何哉？且跖之为人也，心如涌泉，意如飘风，强足以距敌，
辩足以饰非。顺其心则喜，逆其心则怒，易辱人以言。先生必无往。”

孔子不听，颜回为驭，子贡为右，往见盗跖。

\subsection{整数规划}
盗跖乃方休卒徒大山之阳，脍人肝而餔之。孔子下车而前，见谒者曰：“鲁人孔丘，闻将军
高义，敬再拜谒者。”谒者入通。盗跖闻之大怒，目如明星，发上指冠，曰：“此夫鲁国之
巧伪人孔丘非邪？为我告之：尔作言造语，妄称文、武，冠枝木之冠，带死牛之胁，多辞缪
说，不耕而食，不织而衣，摇唇鼓舌，擅生是非，以迷天下之主，使天下学士不反其本，妄
作孝弟，而侥幸于封侯富贵者也。子之罪大极重，疾走归！不然，我将以子肝益昼餔之膳。”

% % !TeX root = ../main.tex

\begin{survey}
\label{cha:survey}

\title{Title of the Survey}
\maketitle

写出至少 5000 外文印刷字符的调研阅读报告或者书面翻译 1-2 篇(不少于 2 万外文印刷符)。

It is impossible to cover in a single chapter every concept of mathematical
programming.\cite{tex} This chapter introduces only the basic concepts and techniques of
mathematical programming such that readers gain an understanding of them
throughout the book\cite{abrahams99tex,salomon1995advanced}.


\section{Single-Objective Programming}
The general form of single-objective programming (SOP) is written
as follows,
\begin{equation*} % 如果附录中的公式不想让它出现在公式索引中，那就请
                             % 用 equation*
\left\{\begin{array}{l}
\max \,\,f(x)\\[0.1 cm]
\mbox{subject to:} \\ [0.1 cm]
\qquad g_j(x)\le 0,\quad j=1,2,\cdots,p
\end{array}\right.
\end{equation*}
which maximizes a real-valued function $f$ of
$x=(x_1,x_2,\cdots,x_n)$ subject to a set of constraints.

\newcommand\Real{\mathbf{R}}
\newtheorem{mpdef}{Definition}[chapter]
\begin{mpdef}
In SOP, we call $x$ a decision vector, and
$x_1,x_2,\cdots,x_n$ decision variables. The function
$f$ is called the objective function. The set
\begin{equation*}
S=\left\{x\in\Real^n\bigm|g_j(x)\le 0,\,j=1,2,\cdots,p\right\}
\end{equation*}
is called the feasible set. An element $x$ in $S$ is called a
feasible solution.
\end{mpdef}

\newtheorem{mpdefop}[mpdef]{Definition}
\begin{mpdefop}
A feasible solution $x^*$ is called the optimal
solution of SOP if and only if
\begin{equation}
f(x^*)\ge f(x)
\end{equation}
for any feasible solution $x$.
\end{mpdefop}

One of the outstanding contributions to mathematical programming was known as
the Kuhn-Tucker conditions\ref{eq:ktc}. In order to introduce them, let us give
some definitions. An inequality constraint $g_j(x)\le 0$ is said to be active at
a point $x^*$ if $g_j(x^*)=0$. A point $x^*$ satisfying $g_j(x^*)\le 0$ is said
to be regular if the gradient vectors $\nabla g_j(x)$ of all active constraints
are linearly independent.

Let $x^*$ be a regular point of the constraints of SOP and assume that all the
functions $f(x)$ and $g_j(x),j=1,2,\cdots,p$ are differentiable. If $x^*$ is a
local optimal solution, then there exist Lagrange multipliers
$\lambda_j,j=1,2,\cdots,p$ such that the following Kuhn-Tucker conditions hold,
\begin{equation}
\label{eq:ktc}
\left\{\begin{array}{l}
    \nabla f(x^*)-\sum\limits_{j=1}^p\lambda_j\nabla g_j(x^*)=0\\[0.3cm]
    \lambda_jg_j(x^*)=0,\quad j=1,2,\cdots,p\\[0.2cm]
    \lambda_j\ge 0,\quad j=1,2,\cdots,p.
\end{array}\right.
\end{equation}
If all the functions $f(x)$ and $g_j(x),j=1,2,\cdots,p$ are convex and
differentiable, and the point $x^*$ satisfies the Kuhn-Tucker conditions
(\ref{eq:ktc}), then it has been proved that the point $x^*$ is a global optimal
solution of SOP.

\subsection{Linear Programming}
\label{sec:lp}

If the functions $f(x),g_j(x),j=1,2,\cdots,p$ are all linear, then SOP is called
a {\em linear programming}.

The feasible set of linear is always convex. A point $x$ is called an extreme
point of convex set $S$ if $x\in S$ and $x$ cannot be expressed as a convex
combination of two points in $S$. It has been shown that the optimal solution to
linear programming corresponds to an extreme point of its feasible set provided
that the feasible set $S$ is bounded. This fact is the basis of the {\em simplex
  algorithm} which was developed by Dantzig as a very efficient method for
solving linear programming.
\begin{table}[ht]
\centering
  \centering
  \caption*{Table~1\hskip1em This is an example for manually numbered table, which
    would not appear in the list of tables}
  \label{tab:badtabular2}
  \begin{tabular}[c]{|m{1.5cm}|c|c|c|c|c|c|}\hline
    \multicolumn{2}{|c|}{Network Topology} & \# of nodes &
    \multicolumn{3}{c|}{\# of clients} & Server \\\hline
    GT-ITM & Waxman Transit-Stub & 600 &
    \multirow{2}{2em}{2\%}&
    \multirow{2}{2em}{10\%}&
    \multirow{2}{2em}{50\%}&
    \multirow{2}{1.2in}{Max. Connectivity}\\\cline{1-3}
    \multicolumn{2}{|c|}{Inet-2.1} & 6000 & & & &\\\hline
    \multirow{2}{1.5cm}{Xue} & Rui  & Ni &\multicolumn{4}{c|}{\multirow{2}*{\thuthesis}}\\\cline{2-3}
    & \multicolumn{2}{c|}{ABCDEF} &\multicolumn{4}{c|}{} \\\hline
\end{tabular}
\end{table}

Roughly speaking, the simplex algorithm examines only the extreme points of the
feasible set, rather than all feasible points. At first, the simplex algorithm
selects an extreme point as the initial point. The successive extreme point is
selected so as to improve the objective function value. The procedure is
repeated until no improvement in objective function value can be made. The last
extreme point is the optimal solution.

\subsection{Nonlinear Programming}

If at least one of the functions $f(x),g_j(x),j=1,2,\cdots,p$ is nonlinear, then
SOP is called a {\em nonlinear programming}.

A large number of classical optimization methods have been developed to treat
special-structural nonlinear programming based on the mathematical theory
concerned with analyzing the structure of problems.
\begin{figure}[h]
  \centering
  \includegraphics{thu-lib-logo.pdf}
  \caption*{Figure~1\quad This is an example for manually numbered figure,
    which would not appear in the list of figures}
  \label{tab:badfigure2}
\end{figure}

Now we consider a nonlinear programming which is confronted solely with
maximizing a real-valued function with domain $\Real^n$.  Whether derivatives are
available or not, the usual strategy is first to select a point in $\Real^n$ which
is thought to be the most likely place where the maximum exists. If there is no
information available on which to base such a selection, a point is chosen at
random. From this first point an attempt is made to construct a sequence of
points, each of which yields an improved objective function value over its
predecessor. The next point to be added to the sequence is chosen by analyzing
the behavior of the function at the previous points. This construction continues
until some termination criterion is met. Methods based upon this strategy are
called {\em ascent methods}, which can be classified as {\em direct methods},
{\em gradient methods}, and {\em Hessian methods} according to the information
about the behavior of objective function $f$. Direct methods require only that
the function can be evaluated at each point. Gradient methods require the
evaluation of first derivatives of $f$. Hessian methods require the evaluation
of second derivatives. In fact, there is no superior method for all
problems. The efficiency of a method is very much dependent upon the objective
function.

\subsection{Integer Programming}

{\em Integer programming} is a special mathematical programming in which all of
the variables are assumed to be only integer values. When there are not only
integer variables but also conventional continuous variables, we call it {\em
  mixed integer programming}. If all the variables are assumed either 0 or 1,
then the problem is termed a {\em zero-one programming}. Although integer
programming can be solved by an {\em exhaustive enumeration} theoretically, it
is impractical to solve realistically sized integer programming problems. The
most successful algorithm so far found to solve integer programming is called
the {\em branch-and-bound enumeration} developed by Balas (1965) and Dakin
(1965). The other technique to integer programming is the {\em cutting plane
  method} developed by Gomory (1959).

\hfill\textit{Uncertain Programming\/}\quad(\textsl{BaoDing Liu, 2006.2})

\bibliographystyle{IEEEtran}
\bibliography{ref/refs,ref/appendix}

\end{survey}
       % 本科生：外文资料的调研阅读报告
% % !TeX root = ../thuthesis-example.tex

\begin{translation}
\label{cha:translation}

\title{书面翻译题目}
\maketitle

\tableofcontents


本科生的外文资料书面翻译。


\section{图表示例}

\subsection{图}

附录中的图片示例（图~\ref{fig:appendix-translation-figure}）。

\begin{figure}
  \centering
  \includegraphics[width=0.6\linewidth]{example-image-a.pdf}
  \caption{附录中的图片示例}
  \label{fig:appendix-translation-figure}
\end{figure}


\subsection{表格}

附录中的表格示例（表~\ref{tab:appendix-translation-table}）。

\begin{table}
  \centering
  \caption{附录中的表格示例}
  \begin{tabular}{ll}
    \toprule
    文件名          & 描述                         \\
    \midrule
    thuthesis.dtx   & 模板的源文件，包括文档和注释 \\
    thuthesis.cls   & 模板文件                     \\
    thuthesis-*.bst & BibTeX 参考文献表样式文件    \\
    thuthesis-*.bbx & BibLaTeX 参考文献表样式文件  \\
    thuthesis-*.cbx & BibLaTeX 引用样式文件        \\
    \bottomrule
  \end{tabular}
  \label{tab:appendix-translation-table}
\end{table}


\section{数学公式}

附录中的数学公式示例（公式~\eqref{eq:appendix-translation-equation}）。
\begin{equation}
  \frac{1}{2 \symup{\pi} \symup{i}} \int_\gamma f = \sum_{k=1}^m n(\gamma; a_k) \mathscr{R}(f; a_k)
  \label{eq:appendix-translation-equation}
\end{equation}


\section{文献引用}

文献引用示例\cite{abrahams99tex}。


% 书面翻译的参考文献
\bibliographystyle{unsrtnat}
\bibliography{ref/appendix}

% 书面翻译对应的原文索引
\begin{translation-index}
  \nocite{salomon1995advanced}
  \bibliographystyle{unsrtnat}
  \bibliography{ref/appendix}
\end{translation-index}

\end{translation}
  % 本科生：外文资料的书面翻译

% 致谢
% !TeX root = ../thuthesis-example.tex

\begin{acknowledgements}
  衷心感谢导师×××教授和物理系××副教授对本人的精心指导。他们的言传身教将使我终生受益。

  在美国麻省理工学院化学系进行九个月的合作研究期间，承蒙 Robert Field 教授热心指导与帮助，不胜感激。

  感谢×××××实验室主任×××教授，以及实验室全体老师和同窗们学的热情帮助和支持！

  本课题承蒙国家自然科学基金资助，特此致谢。
\end{acknowledgements}


% 声明
\statement
% 生成的声明页是否要插入页眉和页脚（默认 empty）
% 仅在需要进行电子签名时，才需要打开这一选项
% 插入的扫描声明页总是会生成页眉（研究生）和页脚，不受这一选项影响
% \statement[page-style=plain]
% 将签字扫描后的声明文件 scan-statement.pdf 替换原始页面
% \statement[file=scan-statement.pdf]

% 个人简历、在学期间完成的相关学术成果
\begin{resume}

  \resumeitem{个人简历}

  xxxx 年 xx 月 xx 日出生于 xx 省 xx 县。
  
  xxxx 年 9 月考入 xx 大学 xx 系 xx 专业，xxxx 年 7 月本科毕业并获得 xx 学士学位。
  
  xxxx 年 9 月免试进入 xx 大学 xx 系攻读 xx 学位至今。

  \researchitem{发表的学术论文} % 发表的和录用的合在一起
% 学位论文写作指南：
% 在学期间发表的学术论文分以下三部分按顺序分别列出，每部分之间空 1
% 行，序号可连续排列
% 1. 已经刊载的学术论文（本人是第一作者，或者导师为第一作者本人是第二作者）
% 2. 尚未刊载，但已经接到正式录用函的学术论文（本人为第一作者，或者
% 导师为第一作者本人是第二作者）。
% 3. 其他学术论文。可列出除上述两种情况以外的其他学术论文，但必须是
% 已经刊载或者收到正式录用函的论文。
  \begin{publications}
  \item Yang Y, Ren T L, Zhang L T, et al. Miniature microphone with silicon-
    based ferroelectric thin films. Integrated Ferroelectrics, 2003,
    52:229-235. (SCI 收录, 检索号:758FZ.)
  \item 杨轶, 张宁欣, 任天令, 等. 硅基铁电微声学器件中薄膜残余应力的研究. 中国机
    械工程, 2005, 16(14):1289-1291. (EI 收录, 检索号:0534931 2907.)
  \item 杨轶, 张宁欣, 任天令, 等. 集成铁电器件中的关键工艺研究. 仪器仪表学报,
    2003, 24(S4):192-193. (EI 源刊.)
  \end{publications}
  \begin{publications}
  \item Yang Y, Ren T L, Zhu Y P, et al. PMUTs for handwriting recognition. In
    press. (已被 Integrated Ferroelectrics 录用. SCI 源刊.)
  \end{publications}
  \begin{publications}[after=] % 最后一个学术论文列表加上 [after=]，避免引入多余空白。
  \item Wu X M, Yang Y, Cai J, et al. Measurements of ferroelectric MEMS
    microphones. Integrated Ferroelectrics, 2005, 69:417-429. (SCI 收录, 检索号
    :896KM.)
  \item 贾泽, 杨轶, 陈兢, 等. 用于压电和电容微麦克风的体硅腐蚀相关研究. 压电与声
    光, 2006, 28(1):117-119. (EI 收录, 检索号:06129773469.)
  \item 伍晓明, 杨轶, 张宁欣, 等. 基于MEMS技术的集成铁电硅微麦克风. 中国集成电路, 
    2003, 53:59-61.
  \end{publications}

  \researchitem{研究成果} % 有就写，没有就删除
  \begin{achievements}
  \item 任天令, 杨轶, 朱一平, 等. 硅基铁电微声学传感器畴极化区域控制和电极连接的
    方法: 中国, CN1602118A. (中国专利公开号.)
  \item Ren T L, Yang Y, Zhu Y P, et al. Piezoelectric micro acoustic sensor
    based on ferroelectric materials: USA, No.11/215, 102. (美国发明专利申请号.)
  \end{achievements}
\end{resume}


% 指导教师/指导小组学术评语
% !TeX root = ../thuthesis-example.tex

\begin{comments}
% \begin{comments}[name = {指导小组评语}]
% \begin{comments}[name = {Comments from Thesis Supervisor}]
% \begin{comments}[name = {Comments from Thesis Supervision Committee}]

  论文提出了……

\end{comments}


% 答辩委员会决议书
% !TeX root = ../thuthesis-example.tex

\begin{resolution}

  论文提出了……

  论文取得的主要创新性成果包括：

  1. ……

  2. ……

  3. ……

  论文工作表明作者在×××××具有×××××知识，具有××××能力，论文××××，答辩××××。

  答辩委员会表决，（×票/一致）同意通过论文答辩，并建议授予×××（姓名）×××（门类）学博士/硕士学位。

\end{resolution}


% 本科生的综合论文训练记录表（扫描版）
% \record{file=scan-record.pdf}

\end{document}
